\documentclass[a4paper,11pt]{article}

% --- Packages ---
\usepackage[utf8]{inputenc} % Character encoding
\usepackage[T1]{fontenc}    % Font encoding
\usepackage[left=0.75in,right=0.75in,top=0.6in,bottom=0.6in]{geometry} % Margins
\usepackage{titlesec}       % Custom section titles
\usepackage{enumitem}       % Custom lists
\usepackage{hyperref}       % Hyperlinks
\usepackage{xcolor}         % Colors

% --- Formatting Setup ---

% Define link colors
\hypersetup{
    colorlinks=true,
    linkcolor=blue,
    filecolor=magenta,      
    urlcolor=teal,
}

% Custom Section Formatting
\titleformat{\section}
{\Large\bfseries\uppercase} % Format
{}                          % Label
{0em}                       % Sep
{}                          % Before
[\titlerule]                % After (horizontal line)

% Custom List Spacing (Tightens up bullet points)
\setlist[itemize]{noitemsep, topsep=0pt, leftmargin=1.5em}

% Custom Command for Experience Headers
% Usage: \experienceheader{Role}{Dates}{Company/Institution}{Location}
\newcommand{\experienceheader}[4]{
    \noindent \textbf{#1} \hfill #2 \\
    \textit{#3} \hfill \textit{#4} \vspace{2pt}
}

% --- Begin Document ---
\begin{document}

% --- Header ---
\begin{center}
    {\Huge \textbf{Abdelhamid FSIAN}} \\ \vspace{5pt}
    Toulouse, France \\
    \href{mailto:abdelhamid-nour-eddine.fsian@u-bourgogne.fr}{abdelhamid-nour-eddine.fsian@u-bourgogne.fr} | \href{mailto:AbdelhamidNour.Fsian@cnes.fr}{AbdelhamidNour.Fsian@cnes.fr} | +33 751 98 67 11\\
    \href{https://www.linkedin.com/in/hamid-fsian-abb669161/}{LinkedIn} | \href{https://github.com/HamidFsian}{GitHub} |
    \href{https://scholar.google.com/citations?user=KpIs5bQAAAAJ&hl=fr&oi=sra}{Google Scholar}
\end{center}

\vspace{10pt}

% --- Professional Summary ---
% --- Professional Summary ---
\section{Profile}
\noindent Computational Imagaging researcher with 4+ years of experience in image reconstruction, deep learning, and material appearance estimation. Developed novel pipelines for joint demosaicking and super-resolution of Spectral Filter Array (SFA) data, and extended them to multi-frame and multi-view video sequences. Strong background in scientific programming, research communication, and deployment of ML models in production environments.

% --- Education ---
\section{Education}

\experienceheader{Ph.D. in Computational Imaging \& Computer Vision}{10/2022 -- 10/2025 }{Université de Bourgogne }{Dijon, France}
\begin{itemize}
    \item \textbf{Thesis:} ``Unified spectral image reconstruction and material appearance estimation from video sequences using Spectral Filter Array (SFA) cameras.''
    \item \textbf{Advisors:} Dr. Jean-Baptiste Thomas, Prof. Pierre Gouton, Prof. Jon Y. Hardeberg
    \item Awarded CNES postdoctoral fellowship to work on lunar rover imaging (RASHID-2 mission).
\end{itemize}
\vspace{8pt}

\experienceheader{M.Sc. in Image Processing \& Computer Vision}{09/2020 -- 09/2022}{Université de Bourgogne}{Dijon, France}
\begin{itemize}
    \item Major de promotion (ranked 1st).
    \item Master's thesis: Joint demosaicking and super-resolution of multispectral video using deep learning.
\end{itemize}
\vspace{8pt}

\experienceheader{M.Sc. in Automation \& Control}{09/2018 -- 07/2020}{University of Tlemcen}{Tlemcen, Algeria}
\begin{itemize}
    \item Specialized in control theory and signal processing.
\end{itemize}

% --- Industrial Experience ---

% --- Experience ---
\section{Experience}

\experienceheader{3D Team Lead}{Dec 2025 -- Present}{Velmeni}{Remote, USA}
\begin{itemize}
    \item Leading the 3D engineering team in developing advanced dental imaging solutions and overseeing the technical roadmap for CBCT processing.
    \item Coordinating cross-functional efforts between ML researchers and backend engineers to accelerate the deployment of diagnostic tools.
    \item Establishing best practices for 3D data versioning and absolute calibration of raw imagery to ensure clinical accuracy.
    \item Managing sprint planning, technical debt reduction, and professional development for a team of specialized engineers.
\end{itemize}
\vspace{4pt}

\experienceheader{Research Scientist}{Oct 2022 -- Apr 2025}{Velmeni}{Remote, USA}
\begin{itemize}
    \item Trained and deployed deep learning models for dental image segmentation (teeth, sinuses, mandibular canal) to assist dentists.
    \item Designed and maintained end-to-end CBCT pipelines (DICOM $\rightarrow$ NIfTI, STL generation, 3D rendering) for treatment planning.
    \item Deployed REST APIs for ML models (FastAPI, Docker, AWS) and integrated them into the clinical web platform.
    \item Improved robustness and latency of 3D verification workflows by refactoring Celery workers and stabilizing background jobs.
    \item Implemented automated evaluation scripts to monitor model performance on real cases (e.g., Dice score, Hausdorff distance).
    \item Standardized internal codebases (logging, configuration, S3 utilities), reducing maintenance overhead and easing onboarding.
    \item Collaborated with dentists and backend engineers to define requirements, validate outputs, and iterate on product features.
    \item Mentored junior engineers and interns on medical image processing and best practices in production ML.
\end{itemize}
\vspace{8pt}


\experienceheader{Postdoctoral Researcher}{Jan 2025 -- Present}{CNES}{Toulouse, France}
\begin{itemize}
    \item Characterization and spectral image reconstruction for the SFA camera onboard the RASHID-2 and 3 lunar rover.
\end{itemize}
\vspace{8pt}


\experienceheader{Co-Owner}{Mar 2024 -- Present}{Spektralion}{Gjøvik, Norway}
\begin{itemize}
    \item Spektralion advises companies on the use of spectral imaging technologies.
\end{itemize}
\vspace{8pt}

\experienceheader{Visiting Researcher}{June 2023 to December 2023}{Norwegian University of Science and Technology (NTNU)}{Gjøvik, Norway}
\begin{itemize}
    \item Conducted collaborative research on spectral imaging and color science within the Department of Computer Science.
\end{itemize}
\vspace{8pt}




\experienceheader{Master's Thesis Co-Supervisor}{Mar 2025 -- Oct 2025}{Université de Bourgogne Europe}{Dijon, France}
\begin{itemize}
    \item Co-supervised a Master’s thesis on applying diffusion models for spatial super-resolution from synthetic spectral data.
    \item Directed the research to publication: \textit{A Comparative Evaluation of Super-Resolution Methods for Spectral Images Using Pretrained RGB Models} (\href{https://doi.org/10.3390/s1010000}{DOI: 10.3390/s1010000}).
\end{itemize}
\vspace{12pt}

\experienceheader{Computer Vision Intern}{Mar 2022 -- Sep 2022}{Actemium Paris Transport}{Paris, France}
\begin{itemize}
    \item Contributed to the development of visual recognition systems for transportation infrastructure.
\end{itemize}


% --- Academic Teaching ---
\section{Academic Teaching}


\experienceheader{Teaching Assistant}{2022 -- 2025}{Université de Bourgogne}{Dijon, France}
\begin{itemize}
    \item Delivered lab sessions and tutorials in image processing, colorimetry, and electronics for undergraduate and master's students.
    \item Supervised research projects on multispectral image stitching and deep learning.
    \item Co-supervised a Master’s thesis that led to the publication:
    \begin{itemize}
        \item \textit{“Evaluating Super-Resolution for Spectral Imaging: A Comparison of Interpolation, CNN, GAN, and Diffusion”}, submitted in 2025.
    \end{itemize}
\end{itemize}

% --- Research Experience ---

\section{Research Experience}

\noindent \textbf{Joint Demosaicking and Super-Resolution for SFA Cameras} | \textit{Python, PyTorch}
\begin{itemize}
    \item Developed a deep learning model for joint reconstruction from raw SFA images (single-frame).
    \item Designed a two-branch architecture with attention-based fusion and pixel-shuffle upsampling.
\end{itemize}

\noindent \textbf{Multi-Frame and Multi-View Spectral Image Fusion} | \textit{Python, OpenCV}
\begin{itemize}
    \item Implemented multi-frame alignment and stitching using homography estimation on selected spectral bands.
    \item Analyzed impact of overlap ratio and keypoint selection on spectral reconstruction quality.
\end{itemize}

\noindent \textbf{SVBRDF Estimation from Multispectral Sequences} | \textit{Python, PyTorch}
\begin{itemize}
    \item Adapted a state-of-the-art SVBRDF estimation pipeline to handle 9-band spectral image sequences.
    \item Replaced RGB backbones with SwinUNet and spectral-aware attention blocks to estimate normals, roughness, and reflectance.
\end{itemize}


%------ LISTS of Publications ---
% --- Publications ---
\section{Publications}

\begin{itemize}
    \item A.N. Fsian, J.-B. Thomas, J.Y. Hardeberg, and P. Gouton, “\textbf{Deep Joint Demosaicking and Super-Resolution for Spectral Filter Array Images},” \textit{IEEE Access}, vol. 13, pp. 16208–16222, 2025. \href{https://doi.org/10.1109/ACCESS.2025.3528753}{DOI:10.1109/ACCESS.2025.3528753}
    
    \item A.N. Fsian, J.-B. Thomas, J.Y. Hardeberg, and P. Gouton, “\textbf{Spectral Reconstruction from RGB Imagery: A Potential Option for Infinite Spectral Data?},” \textit{Sensors}, vol. 24, no. 11, p. 3666, 2024. \href{https://doi.org/10.3390/s24113666}{DOI:10.3390/s24113666}
    
    \item A.N. Fsian, J.-B. Thomas, J.Y. Hardeberg, and P. Gouton, “\textbf{Stitching from Spectral Filter Array Video Sequences},” in \textit{Computational Color Imaging (CCIW 2024)}, LNCS, vol. 15193, Springer, 2025. \href{https://doi.org/10.1007/978-3-031-72845-7_10}{DOI:10.1007/978-3-031-72845-7\_10}
    
    \item A.N. Fsian, J.-B. Thomas, J.Y. Hardeberg, and P. Gouton, “\textbf{Bayesian Multispectral Videos Super Resolution},” in \textit{Proc. 11th European Workshop on Visual Information Processing (EUVIP)}, Gjøvik, Norway, 2023, pp. 1–6. \href{https://doi.org/10.1109/EUVIP58404.2023.10323068}{DOI:10.1109/EUVIP58404.2023.10323068}
    
    \item A.N. Fsian, M.A. Khawaja, M. Kresovic, L. Righetto, and J.Y. Hardeberg, “\textbf{Apple Bruise Detection using Multispectral Imaging},” in \textit{Proc. ICDIP 2025}, Proc. SPIE 13709, 137091Z (2025). \href{https://doi.org/10.1117/12.3075704}{DOI:10.1117/12.3075704}
    
    \item A.N. Fsian, J.-B. Thomas, J.Y. Hardeberg, and P. Gouton, “\textbf{Unified method for image reconstruction and super-resolution of SFA video sequences},” \textit{EURASIP Journal on Image and Video Processing}, under review, 2025.
\end{itemize}


% --- Software Skills ---
\section{Software Skills}
\begin{itemize}
    \item \textbf{Programming:} Python, C++, MATLAB, Bash
    \item \textbf{Deep Learning:} PyTorch, TensorFlow, Keras
    \item \textbf{Image Processing:} OpenCV, SimpleITK, Scikit-Image
    \item \textbf{Tools \& Infra:} Git, Docker, Celery, Traefik, AWS S3, FastAPI
    \item \textbf{OS:} Linux, Windows
\end{itemize}

\section{Languages}
\noindent French (Native), English (Professional), Arabic (Native), Japanese (Basic)

\end{document}